\section {Modelo de ejecución SIMD}
\subsection{¿Por qué una unidad SIMD?}

\begin{frame}{Limitaciones del procesamiento escalar}
 \begin{textblock*}{\textwidth}(0.4cm,1.2cm) \justifying
  \only<1->{
   \begin{tcolorbox}[enhanced jigsaw,rounded corners, drop fuzzy shadow=ShadowColor]
    Durante muchos años los procesadores incrementaron su rendimiento aumentando la frecuencia de reloj y explotando el paralelismo a nivel de instrucciones (\emph{Instruction Level Parallelism, ILP}). Sin embargo, numerosas aplicaciones como procesamiento de imágenes, audio, video y señales digitales presentan una gran cantidad de operaciones idénticas sobre conjuntos extensos de datos.
   \end{tcolorbox}
   }
 \vspace{-0.3cm}
 \only<2->{
  \begin{tcolorbox}[enhanced jigsaw,rounded corners, drop fuzzy shadow=ShadowColor]
   En estos casos, ejecutar una instrucción por cada elemento resulta poco eficiente, aun cuando el procesador sea superescalar y pueda ejecutar varias instrucciones simultáneamente. Aun operaciones elementales como la suma.
  \end{tcolorbox}
 }
 \vspace{-0.6cm}
 \only<3->{
 \begin{center}
  \begin{tabular}{c}
$a_0+b_0$\\
$a_1+b_1$\\
$a_2+b_2$\\
$a_3+b_3$
  \end{tabular}
 \end{center}
 \vspace{-0.3cm}
 \textbf{\textit{\centering Cada suma requiere una instrucción independiente.}}
 }
 \end{textblock*}
\end{frame}

\begin{frame}{Procesamiento SIMD}
 \begin{textblock*}{\textwidth}(0.4cm,1.2cm) \justifying
  \only<1->{
   \begin{tcolorbox}[enhanced jigsaw,rounded corners, drop fuzzy shadow=ShadowColor]
    Las aplicaciones multimedia presentan un elevado paralelismo a nivel de datos (\emph{Data Level Parallelism, DLP}), ya que la misma operación suele aplicarse sobre grandes arreglos o vectores.
   \end{tcolorbox}
   }
 \vspace{-0.3cm}
\only<2->{
  \begin{tcolorbox}[enhanced jigsaw,rounded corners, drop fuzzy shadow=ShadowColor]
   Las arquitecturas SIMD (\emph{Single Instruction Multiple Data}) permiten que una única instrucción opere simultáneamente sobre varios elementos almacenados en un registro ancho, reduciendo el número de instrucciones ejecutadas y aumentando el rendimiento.
  \end{tcolorbox}
  }
%  \vspace{0.5cm}
 \only<3->{
  \begin{center}
   \[
    \left[
     \begin{array}{cccc}
      a_0&a_1&a_2&a_3
     \end{array}
    \right]
     +
    \left[
     \begin{array}{cccc}
      b_0&b_1&b_2&b_3
     \end{array}
    \right]
    \longrightarrow
    \left[
     \begin{array}{cccc}
      c_0&c_1&c_2&c_3
     \end{array}
    \right]
   \]
  \end{center}
  \textbf{\textit{\centering Una única instrucción produce cuatro resultados en paralelo.}}
 }
 \end{textblock*}
\end{frame}

\subsection{Un modelo de paralelización}
\begin{frame}{Single Instruction Multiple Data}
 \begin{textblock*}{0.11\textwidth}(.3cm,1.2cm)
  \includegraphics[width=\textwidth]{SIMD1.pdf}
 \end{textblock*}
 \begin{textblock*}{0.88\textwidth}(.14\textwidth,1.2cm)
  \begin{itemize}[<+(1)->]\justifying
   \item Se trata de un modelo de ejecución capaz de computar una sola operación sobre un conjunto de múltiples datos.
   \item Se refiere a esta técnica como paralelismo a nivel de datos.
   \item Es particularmente útil para procesar audio, video, o imágenes en donde se aplican algoritmos repetitivos sobre sets de datos del mismo formato y que se procesan en conjunto, como por ejemplo en filtros, compresores, codificadores en donde la salida depende de los últimos \textit{n} valores de muestras tomados.
   \item La figura muestra el layout típico de variables memoria para aplicar este modelo
  \end{itemize}
 \end{textblock*}
\end{frame}

\begin{frame}{Como sería la vida sin SIMD?}
 \begin{textblock*}{0.11\textwidth}(.3cm,1.2cm)
  \includegraphics[width=\textwidth]{SIMD1.pdf}
 \end{textblock*}
 \begin{textblock*}{0.89\textwidth}(.11\textwidth,1.2cm)
  \begin{itemize}[<+(1)->]\justifying
   \item En contraposición a SIMD el modelo previo es SISD (\textbf{S}ingle \textbf{I}nstruction \textbf{S}ingle \textbf{D}ata)
   \item Considerando el sector de memoria de la figura, nos proponemos realizar una operación aritmética o lógica sobre las cadenas de datos $a_{n}$ y $b_{n}$, almacenando el resultado en $r_{n}$
   \item Si un procesador no dispone de un modelo arquitectural que le permita implementar paralelismo a nivel de datos, como SIMD,esta operación implica un loop, ya que solo puede procesar un dato por cada instrucción (SISD).
  \end{itemize}
 \end{textblock*}
\end{frame}

\begin{frame}{Single Instruction Single Data}
 \begin{textblock*}{\textwidth}(0.3cm,1.2cm)
  \centering \includegraphics[width=.4\textwidth]{SISD-Layer2.pdf}
 \end{textblock*}
\end{frame}

\begin{frame}{Single Instruction Single Data}
 \begin{textblock*}{\textwidth}(0.3cm,1.2cm)
  \centering \includegraphics[width=.4\textwidth]{SISD-Layer3.pdf}
 \end{textblock*}
\end{frame}

\begin{frame}{Single Instruction Single Data}
 \begin{textblock*}{\textwidth}(0.3cm,1.2cm)
  \centering \includegraphics[width=.4\textwidth]{SISD-Layer4.pdf}
 \end{textblock*}
\end{frame}

\begin{frame}{Single Instruction Single Data}
 \begin{textblock*}{\textwidth}(0.3cm,1.2cm)
  \centering \includegraphics[width=.4\textwidth]{SISD-Layer5.pdf}
 \end{textblock*}
\end{frame}

\begin{frame}{Single Instruction Multiple Data}
 \begin{textblock*}{\textwidth}(0.3cm,1.2cm)
  \begin{itemize}\justifying
   \item Cada Registro se carga en una sola instrucción
   \item La operación se efectúa en la segunda instrucción
  \end{itemize}
  \centering \includegraphics[width=.6\textwidth]{SIMD6.pdf}
 \end{textblock*}
\end{frame}

\begin{frame}{Resumiendo, SIMD vs. SISD}
 \begin{textblock*}{\textwidth}(0.3cm,1.2cm)
  \centering \includegraphics[width=.9\textwidth]{SISD-vs-SIMD.pdf}
 \end{textblock*}
\end{frame}

\begin{frame}{Acerca de SIMD y Procesamiento Vectorial}
 \begin{textblock*}{\textwidth}(0.4cm,1.2cm)
  \begin{tcolorbox}[enhanced jigsaw,rounded corners, drop fuzzy shadow=ShadowColor]
   \begin{itemize}[<+->]\justifying
    \item \textbf{SIMD} aplica una misma instrucción sobre un conjunto fijo de datos empaquetados en un registro ancho. Como ocurre con \textbf{NEON}.
    \item Cada registro con los que opera una instrucción \textbf{SIMD} contiene un grupo fijo de operandos empaquetados. La misma instrucción se ejecuta en paralelo sobre cada operando individual.
    \item El procesamiento vectorial opera sobre vectores de longitud variable mediante registros vectoriales y hardware capaz de recorrer automáticamente todos sus elementos.
    \item Conceptualmente, el programador ``vé'' usa loa instrucción vectorial, por ejemplo \texttt{VADD V1, V2, V3}, es decir, "sumar todos los elementos de V2 y V3 y almacenarlos en V1".
    \item Internamente el hardware ejecuta una secuencia de operaciones elementales sobre cada componente del vector, insertándolas en el pipeline.
   \end{itemize}
  \end{tcolorbox}
 \end{textblock*}
\end{frame}

\begin{frame}{Acerca de SIMD y Procesamiento Vectorial}
 \begin{textblock*}{\textwidth}(0.4cm,1.2cm)
  \begin{tcolorbox}[enhanced jigsaw,rounded corners, drop fuzzy shadow=ShadowColor]
   \begin{itemize}[<+->]\justifying
    \item La diferencia con \textbf{SIMD} es que el programador no tiene que preocuparse por cuántos elementos caben en un registro: la instrucción se aplica al vector completo, aunque el hardware lo procese en varias etapas.
    \item Dicho de otra manera:
    \begin{itemize}[<+->]\justifying
     \item \textbf{SIMD}: una instrucción opera sobre un número fijo de elementos (4, 8, 16, etc.).
     \item \textbf{Vectorial}: una instrucción opera sobre un vector de longitud arbitraria; el hardware se encarga de recorrerlo internamente.
    \end{itemize}
    \item Por eso suele decirse que los procesadores vectoriales tienen vector-length agnostic execution, mientras que \textbf{NEON} es \textbf{SIMD} de ancho fijo: \SI{128}{\bit}.
    \item \textbf{SIMD} puede interpretarse como un caso particular de procesamiento vectorial, en el que la longitud del vector coincide con el ancho fijo del registro.
    \item ARMv8 evoluciona de NEON a SVE adoptando un modelo de procesamiento vectorial de longitud variable  de hasta \SI{1024}{\bit} en pasos de \SI{128}{\bit}.
   \end{itemize}
  \end{tcolorbox}
 \end{textblock*}
\end{frame}
